\scanpage{64}\mspage{4.}
% Long bracket beside the initial setup.
\diagram{\begin{tikzcd}[ampersand replacement=\&,column sep=small,row sep=small]
\& X'\ \text{normalisé}\arrow[d] \& {}\\
Y\subset \& X\text{ -- }\underset{\text{\unsure{réduit|réduite}}}{\text{\uline{\unsure{complète|complet}}}} \& U=X-Y
\end{tikzcd}}

$Y_1$ -- $Y_r$ les composantes de $Y$ de codim $1$ ds $X$

$Y'_{ij}$ les images inverses des $Y_i$ ds $X'$

\deleted{\unsure{$\pi'_1$}\ \uline{Th} (?) $\pi_1(U)$ \inserted{le quotient de $\pi_1(U)$ par le} le noyau de $\pi_1(U)\to\pi_1(X)$
est \unsure{son sous-groupe|un sous-groupe} de} Considérons les rev. étales de $X-Y$ qui
\deleted{\unsure{sont de t. fini|se définissent|sont classifiés}} donnent des rev. \uuline{doux} des $\underline{O}_{y_{ij}}$,
ils sont classifiés par un groupe de type galoisien $\pi$.

\uline{\unsure{Thm.|Th.}} Le groupe $\pi$ est top. de génération finie.

\uline{Cor. 1} Soit $X$ complète, $Y$ de codim $\geq 2$, \inserted{$U=X-Y$.} Alors $\pi_1(\unsure{U|X}\deleted{\unsure{-Y}})$
est top. de présentation finie

\uline{Cor 2} Soit $X$ \deleted{de type} rég. Alors $\pi'_1(X)$ est top. de
type fini
% Two short horizontal ruling strokes below the text. Faint offset strokes across the lower half.
% archive sheet number: 63
